\chapter{Reasons and Rationality}
\label{sec:p11-reasons-and-rationality}

Reasons, successful action, and coherent attitudes answer different questions.
Someone may act for an intelligible reason, fail to achieve what there is reason
to achieve, and retain the same beliefs and intentions throughout. Conversely,
a fortunate outcome need not repair an incoherent intention. This chapter
develops John Broome's investigation of whether rationality can nevertheless
be identified with, or explained entirely through, correct responsiveness to
normative reasons. Its source is his complete chapter~2.1 in
\emph{The Handbook of Rationality}, including its qualifications and notes
\parencite[pp.~129--136]{Broome2021ReasonsRationality}.

The discussion is philosophical argument. Its examples stipulate circumstances
and invite distinctions between assessments; they do not report experiments
on how people actually reason. The distinction between a benchmark and
observed performance in Definition~\ref{def:p19-benchmark} therefore matters
here too. As in Section~\ref{sec:p11-method-boundary}, introducing precise
terms does not convert a disputed philosophical premise into a logical theorem.
Broome's account of rationality and his objections to rival accounts are
identified as such throughout.

\section{Two normative questions and two roles for reasons}
\label{sec:p11-reasons-normativity}

\begin{hbdefinition}{def:p11-reasons-normative-senses}{Correctness and normative authority}
A requirement supplies a \emph{standard of correctness} when conduct or a state
can conform to it or violate it. A further question is whether someone
\emph{ought} to conform, or has a \emph{normative reason} to conform. Calling a
requirement normative in the first sense does not by itself settle the second.
Here an \emph{ought} concerns what the person's combined normative reasons
require, within the account of reasons being considered.
\end{hbdefinition}

Rationality supplies standards in the first sense: examples include avoiding
contradictory beliefs and appropriately connecting intentions about ends and
means. Broome asks whether rationality also has authority in the second sense.
His contrast is with fashion. A convention about what clothing counts as
fashionable can be sufficiently definite to classify an outfit without thereby
establishing that anyone has a reason to follow it. His 2021 example concerns
bellbottoms; its argumentative role does not depend on what happens to be
fashionable at the time of reading. The existence of a standard and its
authority require different arguments
\parencite[p.~129]{Broome2021ReasonsRationality}.

\begin{hbdefinition}{def:p11-reasons-roles}{Motivating and normative reasons}
A \emph{motivating reason} explains, or helps explain, why a person does
something in the relevant motivational way. A \emph{normative reason}
explains, or helps explain, why the person ought to do it. These are distinct
roles a consideration may play. They need not be occupied by disjoint
considerations: the same consideration can both count in favor of an action
and motivate it. An explanation of an event need not be a motivating reason;
an equipment failure can explain a delay without motivating anyone.
\end{hbdefinition}

Broome presents this distinction as an achievement of twentieth-century
philosophy of normativity, citing Smith's 1994 discussion and his own fuller
definitions. That attribution locates the distinction in his argument;
it does not supply those separate works' complete accounts.

The claim that normative reasons are philosophically fundamental became
influential in the work Broome associates with Nagel, Parfit, and Scanlon.
Its relevance here is a proposed reduction: if rationality consisted entirely
in responding correctly to normative reasons, an account of those reasons and
of correct response might supply the whole account of rationality. Merely
observing that rationality and reasons are connected does not establish that
reduction. Nor does identifying someone's motivation establish that the
motivation was normatively adequate
\parencite[p.~129 and p.~135, n.~1]{Broome2021ReasonsRationality}.

\section{Which meaning of rationality is at issue?}
\label{sec:p11-reasons-meaning}

Broome's historical discussion separates the common linguistic root from the
philosophical question. As he reports from historical dictionaries and Day's
study, English \emph{reason} entered through French around the Norman period.
An early witness, the \emph{Ancrene Riwle} manuscript of about 1225, uses the
count noun for explanations and reasons, while also using a non-count sense
for a mental faculty. One passage gives reasons why a person ought to withdraw
from the world, including attaining heaven; another represents the faculty
as a woman in a parable and glosses it as \emph{wit's skill}.
The point is that the faculty sense and the countable-consideration senses were
already distinguishable. An explanation of a delay, a consideration that
motivates, and a consideration that justifies need not perform the same role.

On the lexical history Broome cites, \emph{rational} first appears in English
in 1398 for possession of the faculty of reason, initially predicated of
people, creatures, souls, and minds. \emph{Rationality}, attested in 1627,
names that property. Shakespeare's presentation of reason as guiding the will
in \emph{A Midsummer Night's Dream}, Act~2, Scene~2, illustrates the faculty
usage. These are the source's historical claims,
used to explain a choice of meaning; they are not an independent etymological
investigation undertaken here. Common ancestry in Latin \emph{ratio} does not
by itself prove that a present philosophical concept of rationality must be
defined using normative reasons
\parencite[pp.~130--131]{Broome2021ReasonsRationality}.

\begin{hbdefinition}{def:p11-reasons-structural}{Broome's mental and structural conception}
On the conception used in this chapter's reconstruction of Broome, rationality
concerns both a faculty of the mind and the coherent organization of mental
states that could result from exercising it. \emph{Structural rationality}
emphasizes relations among attitudes, such as beliefs and intentions.
Possessing an ability to connect means with ends is different from actually
having appropriately connected intentions on a particular occasion.

\emph{Mental supervenience}, as Broome uses it, says that fixing all the other
relevant mental properties across possible situations fixes the person's
degree of rationality. This is his commitment about the concept at issue,
not a measured psychological law or a definition binding every use of the
word rationality. Broome cites Wedgwood's 2002 discussion in connection with
this dependence claim.
\end{hbdefinition}

The interpretation of the mental base matters. If knowledge, including its
factive relation to the world, counts as a mental state, two people with the
same intrinsic beliefs can differ in their mental states so understood.
Arguments that hold the mind fixed must say what they hold fixed. A stipulated
equality of beliefs is not automatically equality on every externalist theory
of mind.

Broome also recognizes derived uses. From 1598, in his reported chronology, the
adjective is applied to things that do not themselves have minds: actions,
beliefs, and plans for cities. Beliefs are mental states, but they are not
themselves subjects possessing minds. A city plan can be called rational
because rational designers could produce it; an action can be assessed by
whether performing it would leave its agent no less rational than not
performing it. Such extensions do not force the basic property of the
person to become responsiveness to objective reasons
\parencite[pp.~130--131]{Broome2021ReasonsRationality}.

Some philosophical accounts, including the Kolodny--Brunero discussion
Broome cites, distinguish structural rationality from
\emph{substantive rationality}, using the latter for doing what one has reason
or ought to do. Broome accepts roughly the first characterization and disputes
the second as an account of the ordinary central meaning. His directional
case separates an objective verdict from the agent's perspective: the person
ought to turn left, but misleading evidence supports a firm judgment that
turning right is required. He additionally stipulates that no other attitude
favors turning left. Under that stipulation, doing the objectively required
thing is irrational on Broome's judgment of the mental sense.

The additional stipulation is essential. Arpaly's discussion of Huckleberry
Finn, as Broome reports it, permits an action to be rational through its
coherence with other attitudes even when it conflicts with the agent's
express moral judgment: Huck conceals the escaped slave Jim despite believing
he ought not to do so. The directional example is designed to exclude that
explanation. Broome's appeal to ordinary usage is a conceptual judgment,
not a survey showing unanimous public acceptance. His note tentatively
traces the substantive label to Weber through Kalberg's discussion. He
associates structural usage with Scanlon and Wallace and notes a different
usage in chapter~10.5 of the source handbook. None of these labels should
be treated as a shared definition merely because their wording matches
\parencite[p.~131 and p.~135, nn.~3--5]{Broome2021ReasonsRationality}.

\section{Can rationality be identical with correct response to actual reasons?}
\label{sec:p11-reasons-identity}

\begin{hbdefinition}{def:p11-reasons-response}{Combined requirement, compliance, and basis}
Let \(F\) stand schematically for something a person might do or for an
attitude or arrangement of attitudes the person might have. A reason in favor
of \(F\) need not require \(F\), because other reasons may oppose, override,
or cancel it. In the following discussion, saying that the combined reasons
require \(F\) means that the person ought to \(F\). No additive numerical
theory of weights is assumed, and the discussion assumes that the combined
reasons do not require both \(F\) and not-\(F\).

The person \emph{complies} when the combined reasons require \(F\) and the
person actually \(F\)s. A stronger account of \emph{correct response} can also
require that the person comply because of those reasons. Compliance,
appropriate motivation, intention, and successful performance must therefore
be distinguished before evaluating a reduction.
\end{hbdefinition}

The strongest proposal identifies rationality with correct response to
normative reasons. Broome identifies Kiesewetter's treatment as a sustained
defense and also names Gibbard and Lord; he takes Lord's formulation in
terms of doing what one ought to do to coincide with the relevant proposal.
Broome's initial objection compares their dependence on
the mind. Rationality, on his chosen conception, supervenes on mental
properties; correct response to reasons apparently does not. If the latter
can differ while the relevant mental properties remain fixed, the two
properties cannot be identical.

For compliance, it is useful to distinguish what the reasons require from
whether the person does it. If both are fixed by the mind, compliance is
fixed too. The converse is not a general logical truth: a compound property
can in principle remain fixed while one component varies. Broome explicitly
allows this possibility and doubts that it supplies the required explanation
in the cases at issue. His argument must accordingly be evaluated through
its examples and its account of mind, rather than through an invalid inference
from the grammatical decomposition alone
\parencite[pp.~131--132 and p.~135, nn.~6--8]{Broome2021ReasonsRationality}.

\subsection{The external circumstances that supply reasons}

On an ordinary interpretation, gathering clouds can be a reason to expect
rain, and a child's injury can be a reason to take the child to hospital.
The relevant facts reach beyond someone's present mental state. A strongly
subjective account could instead make what one ought to do depend entirely
on one's own credences and value judgments, for example through a subjective
expected-value criterion. Broome finds the cost of removing the apparent
normative role of worldly facts unattractive. That judgment supplies a
premise against such subjectivism; it is not a theorem against every possible
subjective theory.

A more specific reply, which he associates with Kiesewetter, restricts reasons
for a person to facts available in that person's evidence. The evidence
together with other mental features is then supposed to fix what ought to be
done. Williamson's identification of evidence with knowledge, combined with
the view that knowledge is a mental state, would make the evidence part of
the mental base. Broome objects that this imports the relevant worldly facts
into the mind instead of removing their role. A defender can contest that
assessment, but must specify the enlarged mental base. One cannot accept the
externalist reply and then test it using counterparts matched only in
nonfactive belief
\parencite[p.~132]{Broome2021ReasonsRationality}.

\subsection{The external success of the response}

A second difficulty survives even when the reasons themselves are described
mentally. Broome's agent wants to avoid financial ruin and believes that
insurance is necessary. The agent fills out a form, skims the contract, and
pays. In one situation this obtains insurance. In another, an unnoticed
condition excludes the agent's cedar-shingle roof: the contract requires
one of several other roofing materials. The agent never makes a claim and
never discovers the problem. The relevant mental profile is stipulated to
remain the same, while actual insurance coverage differs.

The example concerns equality of rationality between the situations; it
does not require that either agent be perfectly rational. If the normative
requirement was to secure actual insurance, compliance can differ despite the
stipulated mental equality. Replacing that requirement by one to intend to
insure, or merely to seem to have insured, avoids this performance problem
only by changing the normative account. Broome's objection is that the
agent's interest concerns actual protection, with intention and appearance
valuable insofar as they help obtain it. His note discusses a suggested
appeal to Prichard without claiming to give a complete examination of
Prichard's position
\parencite[p.~132 and p.~135, n.~10]{Broome2021ReasonsRationality}.

\subsection{The further distinction from morality}

Even overcoming both dependence objections would not establish identity.
Broome adds that responding appropriately to moral reasons is an exercise
of morality, and does not thereby amount simply to rationality. Restricting
responses to mental states does not remove this issue: a moral assessment of
racist beliefs or malicious intentions is still a moral assessment. This
does not imply that moral and rational assessments never agree, or that one
mental state cannot possess both kinds of merit. It disputes the identification
of the kinds of assessment. The three objections have different roles:
the first concerns what fixes the normative verdict, the second what fixes
successful compliance, and the third what property the response exemplifies
\parencite[pp.~132--133]{Broome2021ReasonsRationality}.

\section{Would correct response at least guarantee rationality?}
\label{sec:p11-reasons-entailment}

A weaker proposal says that responding correctly to reasons entails being
rational, even if the properties are not identical. An entailment alone
would still leave open whether rationality has an independent explanatory
role. Broome considers two routes to the stronger explanatory claim.

\subsection{Structural reasons and competing reasons}

The first route posits normative reasons to maintain structural coherence,
including consistency and suitable means--end connections. Responses to
these reasons can be entirely mental, so the performance objection to
insurance is avoided. The difficulty is that correct response concerns
all relevant reasons, not a selected group labeled structural.

Broome's demon case stipulates that catastrophe will occur unless a person
holds contradictory beliefs. If the reason to avert catastrophe defeats
the structural reason for consistency, correct response to the combined
reasons would produce the very contradiction that excludes full rationality
on his conception. The contrast is between a requirement whose violation
excludes full rationality and a reason that may be defeated in the overall
normative assessment. It does not assume a numerical scale for comparing
those reasons. A critic may reject the case's normative stipulation, or
deny that the purported threat supplies the appropriate kind of reason for
belief. Those are substantive challenges to its premises; the thought
experiment does not settle them simply by being describable
\parencite[p.~133]{Broome2021ReasonsRationality}.

\subsection{The proposal that coherence has no independent requirements}

The second route is the position Broome calls the \emph{myth theory},
associated in his discussion with Kolodny and Raz. It need not deny that
mental states have structural relationships. It denies that structural
rationality supplies independently interesting prescriptions: responding
properly to the reasons supplied by the world is supposed to produce the
needed coherence automatically.

Two meanings of \emph{requires} must be kept apart here. A property can
require something merely as a necessary condition: having little hair
is necessary for baldness. Alternatively, a requirement can be a prescription
issued by an evaluative standard. The dispute concerns whether rationality
has requirements in the latter sense. It is not answered by observing that
a definition of coherence has necessary conditions
\parencite[pp.~133--134]{Broome2021ReasonsRationality}.

Broome's first reply uses a trip to Paris. Suppose the normative reasons
permit going and permit not going. That alone does not imply permission
to form both an intention to go and an intention not to go. The source
adds the premise that the world supplies no reason against the combination,
and suggests that preparing for each possibility could be useful. Under
those extra assumptions, a person can respond acceptably to world-facing
reasons while having incompatible intentions. Broome treats such practical
inconsistency as excluding full rationality. He reports Kolodny's response
as giving up the strict consistency requirement and regards that price as
too high. The disagreement is over that requirement and the normative
premises; it is not resolved by confusing two permissible options with a
permissible conjunction of attitudes.

His second reply concerns reasoning itself. Finding out which means serve
an end can call for theoretical reasoning; forming an intention to take a
means can call for instrumental reasoning. These involve transitions from
premise attitudes to a conclusion attitude. Broome argues that the
correctness of such transitions depends on rational relationships between
the attitudes, described more fully elsewhere through permissions concerning
what an attitude may be based on. The same inferential relation can apply
when a premise is false or an end is bad. Thus the correctness of a transition
cannot simply be identified with the truth or moral worth of its starting
points.

On this reply, even if ideal responsiveness to reasons always brought
coherence with it, understanding that responsiveness would still require an
account of reasoning. Rationality would retain an explanatory role. This
is an argument for its indispensability, not a derivation of the truth of
every starting belief or the acceptability of every end. The formal
consequence relation in Definition~\ref{def:p11-model-consequence} similarly
distinguishes preservation from the separate question whether the premises
are satisfied, although Broome's broader account of transitions among
attitudes is not identified with that particular formal relation
\parencite[p.~134 and p.~135, n.~13]{Broome2021ReasonsRationality}.

\section{Responding to beliefs about reasons}
\label{sec:p11-reasons-believed}

Moving from actual reasons to beliefs about reasons creates several different
proposals. They should not be treated as interchangeable repairs.

\begin{hbdefinition}{def:p11-reasons-belief-variants}{Three belief-based proposals}
The first proposal requires doing what would be required by the reasons
one believes to exist. The second requires doing what one believes one's
combined reasons require. The third restricts the response in the second
proposal to an intention: intend to do what one believes one's combined
reasons require. A belief that certain considerations exist, a belief about
their combined verdict, an intention to comply, and actual compliance are
four distinguishable items.
\end{hbdefinition}

The first proposal, associated here with Parfit, leaves open how the supposed
reasons in fact combine. That normative relation can depend on matters beyond
the agent's belief, so Broome's first dependence objection remains. The
second, associated with Kolodny, makes the relevant verdict a belief of
the person and thereby avoids that objection. If it still demands a
successful worldly action, however, the insurance objection remains.
Restricting the demanded response to a mental state avoids that second
objection. The intention version makes the relevant relation one between
a judgment about what ought to be done and an intention
\parencite[p.~134]{Broome2021ReasonsRationality}.

Avoiding the two objections is not enough to define all rationality.
Someone may have contradictory beliefs or intentions without believing that
there is anything wrong with the combination. On Broome's conception,
failure to recognize the defect does not abolish it. He calls this aspect
of rational requirements \emph{strict liability}: a violation need not be
accompanied by awareness of the requirement. The expression here makes
no claim about legal liability, blameworthiness, or the psychological ease
of detecting every inconsistency.

\begin{hbdefinition}{def:p11-reasons-enkrasia}{The limited connection: enkrasia}
Broome gives the following rough description of \emph{enkrasia}: rationality
requires a person to intend to do something conditional on the person's
believing that their combined reasons require doing it. The corresponding
failure, called \emph{akrasia} in this formulation, is believing that one ought to do it while
lacking the intention to do it. This is one connection among attitudes,
not a complete definition of rationality.
\end{hbdefinition}

The source explicitly calls its formulation rough and refers to a more
complicated treatment in Broome's 2013 book. No fully qualified formal
enkrasia principle is supplied by the short formulation alone. In particular,
it does not prove that the judgment about reasons is true, that actual
reasons require the action, that the action will succeed, or that every
rational requirement itself supplies a normative reason for compliance.
Nor does detecting a mismatch specify, without further premises, whether
the agent should change the judgment or change the intention.

The historical association with weakness of will reaches back to Aristotle
and is developed in Davidson's discussion, as the source notes. Those
separate works are not reconstructed here. In his concluding section~6,
Broome rejects both identities: rationality with correct response to actual
reasons, and rationality with correct response to beliefs about reasons.
The relation between an ought judgment and intention supplies a limited
surviving connection. His preceding discussion also argues that even a
reasons-response entailment would not eliminate rationality's explanatory
role. These conclusions do not answer every wider question about the
authority of rationality
\parencite[p.~134]{Broome2021ReasonsRationality}.

\section{Using the distinctions in a decision problem}
\label{sec:p11-reasons-diagnostic}

The finite-preference framework in Part~\ref{part:p17} supplies a useful comparison without
settling the philosophical dispute. On a finite set, a complete, transitive ordering has an
ordinal numerical representation by
Theorem~\ref{thm:p17-finite-representation}. That result concerns the
representability of comparisons. It does not establish that the ordering
expresses correct normative reasons, that the agent will follow it, or that
all the agent's attitudes are coherent. Definition~\ref{def:p17-ordinal-interpretation-boundary}
also prevents interpreting rank differences as strengths of reasons.
An appeal to expected value would require the additional probability and
value assumptions of an appropriate decision model.

\begin{hbmethod}{def:p11-reasons-diagnostic}{Diagnosing a reasons--rationality claim}
For a particular agent and occasion, first specify the object of assessment:
an action, a belief, an intention, a combination of attitudes, or successful
achievement of an end. State what counts as its occurrence or success.

Next identify the sense of normativity being used. Is the claim about
conformity to a stated standard, or about a reason or an ought to conform?
Record separately the actual considerations, the considerations the agent
believes exist, the agent's judgment about their combined force, the
intention, and the achieved result. Leave unknown entries unknown.

State the proposed way in which reasons combine. If the claim relies on
weights, priorities, or cancellation, give that premise; a preference rank
alone supplies none of these. Identify whether the proposed conclusion is
identity of properties, sufficiency of one for another, or satisfaction
of one selected coherence requirement.

For a comparison across possible situations, specify the mental base and
the external facts held fixed or changed. Check each proposed inference
against the appropriate distinctions: permission for either option is not
automatically permission for incompatible joint intentions, and an intention
is not automatically a successful action.

Conclude separately about motivation, the stated coherence requirements,
the normative verdict under the premises, and performance. If a needed
normative or mental premise is disputed, identify it rather than treating
the procedure as a calculation of all reasons.
\end{hbmethod}

This is an editorial procedure for using the chapter's distinctions.
It is not an algorithm proposed by Broome or a decision procedure for
normative truth. It can expose missing premises even when the final ought
remains unsettled.

\begin{hbexample}{def:p11-reasons-deposit-case}{A promised deposit and an unnoticed failure}
Mira has promised to deposit a report before a deadline. Stipulate that her
combined normative reasons require an actual deposit. She believes this,
intends to make the deposit, and intends to use an available channel.
Her available alternatives are the portal, email, and waiting. For this
example only, define a preference ordering by assigning rank \(1\) to the
portal and email and rank \(0\) to waiting, with weak preference meaning
that the assigned rank is at least as large.

The rank comparison is complete because any two real ranks are comparable,
and transitive because their numerical order is transitive. The maximal
choice set for the three-option menu is exactly the portal and email:
each has rank at least that of every alternative, while waiting has rank
below both. This is the construction used in
Theorem~\ref{thm:p17-maximal-choice}. Replacing their ranks by \(5,5,-2\)
preserves every comparison: the transformation taking an old rank to
seven times that rank minus two is strictly increasing. It does not show
that the reason against waiting has acquired a different strength.
Choosing the portal from the two eligible channels also does not require
an intention both to deposit and not to deposit.

Mira selects the portal and uploads the report. Compare two stipulated
server situations. In the first, the report is accepted; in the second,
it is rejected without a visible notice. Mira never learns of the difference.
Hold fixed her beliefs, intentions, and all the other mental properties
included in the conception being used. She complies with the stipulated requirement of actual deposit
in the first situation and fails to comply in the second. The local
judgment--intention relation and her intended means are unchanged. Under
Broome's supervenience commitment, her degree of rationality is therefore
unchanged, although this does not establish that her entire mental life
is fully rational.

Her belief about the promise and deadline can explain the upload as a
motivating reason. The promise and circumstances, under the example's
normative stipulation, help justify the requirement to deposit. Those
explanatory roles must be distinguished even if the consideration that
motivates her also counts in favor of the act.

Finally change the mental case: she believes she ought to deposit, has
no intention to do so, and intends to wait. Suppose an automatic archival
process nevertheless deposits the report. Successful performance of the
required outcome does not supply the missing intention. The rough
enkrasia description diagnoses a mismatch in this variant, subject to
its stated limited scope. The example does not infer a psychological
frequency, a legal verdict about promises, or a uniquely prescribed repair.
\end{hbexample}

\section{Scope of the conclusion}
\label{sec:p11-reasons-boundary}

The account makes several assessments available without combining them into
one score: the explanatory role of a motive, the force of normative reasons
under specified premises, structural relations among attitudes, and actual
success. Broome's case against the reductions depends on his conception of
rationality, the relevant interpretation of mind, and the normative premises
of his examples. Its opponents can contest those premises; they cannot
answer an objection about successful performance merely by redescribing
performance as intention without acknowledging the changed proposal.

The source contains conceptual analysis, arguments, historical remarks,
and hypothetical cases, with no participant sample or empirical estimate.
Applying its distinctions to observed reasoning would require the task,
measure, population, and evidential qualifications specified in
Definitions~\ref{def:p19-empirical-question} and~\ref{def:p19-evidence-record}
and Section~\ref{sec:p19-warranted-conclusion}. Likewise, the finite
representation results in Part~\ref{part:p17} supply exact mathematical conclusions
under their ordering assumptions; they do not select among all the
philosophical accounts discussed here.

This treatment covers the selected 2021 chapter, including its thirteen
notes. Its fuller treatments of basing permissions, precise enkrasia
conditions, and several replies are referred to other works, particularly
Broome's 2013 book and his distinct 2021 essay on rationality and normativity.
Those are separate source presentations. The present chapter supplies
neither their unprinted details nor a claim that the whole philosophical
theory of rationality has been completed
\parencite[pp.~135--136]{Broome2021ReasonsRationality}.
